% ============================================================
\chapter{Kakutani's Theorem and Correspondences}
\label{ch:kakutani}
% ============================================================

In 1941, Shizuo Kakutani, a young Japanese mathematician trained at Osaka University, published a result that appeared to be a mere technical extension of Brouwer's theorem: instead of requiring a function to have a fixed point, he showed that a \emph{correspondence} --- a map that assigns to each point not a single value but an \emph{entire set} of values --- also possesses a fixed point, provided the images are convex and the correspondence is upper semicontinuous. The result could have remained confined to topological circles. But nine years later, John Nash used it as the centrepiece of his proof that non-cooperative games possess equilibria, work that would earn him the Nobel Prize in Economics in 1994. Kakutani's theorem thus became one of the pillars of game theory and mathematical economics.

\begin{intuition}
Kakutani's theorem generalizes Brouwer's theorem to set-valued maps (correspondences).
Instead of a function $f\colon C \to C$, one considers a correspondence
$F\colon C \rightrightarrows C$ that associates to each point a \emph{set} of values.
Kakutani shows that such a correspondence has a fixed point (a point $x^*$ with
$x^* \in F(x^*)$) under convexity and semicontinuity hypotheses. This result is
fundamental in game theory, where it is used to prove the existence of Nash equilibria.
\end{intuition}

% ------------------------------------------------------------
\section{Correspondences and Semicontinuity}
% ------------------------------------------------------------

\begin{definition}[Correspondence (set-valued map)]
\index{correspondence}
Let $X$ and $Y$ be topological spaces. A \emph{correspondence} (or set-valued map)
$F\colon X \rightrightarrows Y$ is a map $F\colon X \to \mathcal{P}(Y)$, i.e., for each
$x \in X$, $F(x)$ is a subset of $Y$. The \emph{graph} of $F$ is
$\mathrm{Gr}(F) = \{(x,y) \in X \times Y : y \in F(x)\}$.
\end{definition}

\begin{definition}[Upper hemicontinuity (u.h.c.)]
\index{upper hemicontinuity}
A correspondence $F\colon X \rightrightarrows Y$ is \emph{upper hemicontinuous} (or
\emph{upper semicontinuous}) if for every open set $V \subseteq Y$, the set
\[
  F^+(V) = \{x \in X : F(x) \subseteq V\}
\]
is open in $X$.
\end{definition}

\begin{definition}[Lower hemicontinuity (l.h.c.)]
\index{lower hemicontinuity}
A correspondence $F\colon X \rightrightarrows Y$ is \emph{lower hemicontinuous} if for
every open set $V \subseteq Y$, the set
\[
  F^-(V) = \{x \in X : F(x) \cap V \neq \emptyset\}
\]
is open in $X$.
\end{definition}

\begin{proposition}[Closed graph characterization]
\label{prop:closed-graph}
Let $F\colon X \rightrightarrows Y$ be a compact-valued correspondence and $Y$ a compact
Hausdorff space. Then $F$ is u.h.c.\ if and only if its graph $\mathrm{Gr}(F)$ is closed
in $X \times Y$.
\end{proposition}

\begin{proof}
$(\Rightarrow)$: Let $(x_\alpha, y_\alpha)$ be a net in $\mathrm{Gr}(F)$ converging to
$(x, y)$. Suppose $y \notin F(x)$. Since $F(x)$ is compact and $Y$ is Hausdorff, there
exists an open set $V$ containing $F(x)$ with $y \notin \overline{V}$. By u.h.c.,
$F^+(V)$ is a neighborhood of $x$, so $F(x_\alpha) \subseteq V$ for large $\alpha$,
hence $y_\alpha \in V$. In the limit, $y \in \overline{V}$, contradiction.

$(\Leftarrow)$: Let $V$ be open with $F(x) \subseteq V$. If $F^+(V)$ were not a
neighborhood of $x$, there would exist a net $x_\alpha \to x$ with
$F(x_\alpha) \not\subseteq V$. One could find $y_\alpha \in F(x_\alpha) \setminus V$.
By compactness of $Y$, one can extract a subnet converging to $y \in Y \setminus V$.
Then $(x,y) \in \mathrm{Gr}(F)$ since the graph is closed, so $y \in F(x) \subseteq V$,
contradiction.
\end{proof}

\begin{remark}
For a single-valued function $f\colon X \to Y$ (viewed as a correspondence
$F(x) = \{f(x)\}$), upper hemicontinuity coincides with ordinary continuity.
\end{remark}

% ------------------------------------------------------------
\section{Kakutani's Theorem}
% ------------------------------------------------------------

\begin{theorem}[Kakutani, 1941]
\label{thm:kakutani}
\index{Kakutani's theorem}
Let $C$ be a nonempty compact convex subset of $\R^n$ and
$F\colon C \rightrightarrows C$ a correspondence such that:
\begin{enumerate}
  \item for every $x \in C$, $F(x)$ is nonempty and convex;
  \item $F$ is upper hemicontinuous (closed graph).
\end{enumerate}
Then $F$ has a fixed point: there exists $x^* \in C$ with $x^* \in F(x^*)$.
\end{theorem}

\begin{proof}
For each $n \in \N^*$, consider a simplicial triangulation of $C$ with mesh less than
$1/n$, with vertices $v_1^n, \ldots, v_{k_n}^n$. For each vertex $v_i^n$, choose
$w_i^n \in F(v_i^n)$. Define the piecewise affine map $f_n\colon C \to C$ by
$f_n(v_i^n) = w_i^n$ and affine extension on each simplex.

Since $C$ is convex compact and $f_n$ is continuous, by Brouwer's theorem there exists
$x_n \in C$ with $f_n(x_n) = x_n$.

By compactness, $(x_n)$ has a convergent subsequence $x_{n_k} \to x^*$. We show that
$x^* \in F(x^*)$. For every $\varepsilon > 0$, for large $k$, the mesh of the
triangulation is less than $\varepsilon$, and $x_{n_k}$ belongs to a simplex whose
vertices $v_i^{n_k}$ satisfy $\norm{v_i^{n_k} - x_{n_k}} < \varepsilon$. Since
$f_{n_k}(x_{n_k}) = x_{n_k}$ is a convex combination of the
$w_i^{n_k} \in F(v_i^{n_k})$, and since $\mathrm{Gr}(F)$ is closed
(Proposition~\ref{prop:closed-graph}), we conclude $x^* \in F(x^*)$.
\end{proof}

\begin{attention}
The convexity hypothesis on $F(x)$ is essential. Consider $C = [0,1]$ and
$F(x) = \{0, 1\} \setminus \{x\}$: this closed-graph correspondence has no fixed point
because $F(x)$ is not convex.
\end{attention}

% ------------------------------------------------------------
\section{The Glicksberg--Fan Extension}
% ------------------------------------------------------------

\begin{theorem}[Glicksberg--Fan]
\label{thm:glicksberg-fan}
Let $E$ be a locally convex topological vector space and $C \subseteq E$ a nonempty
compact convex set. Let $F\colon C \rightrightarrows C$ be a correspondence with nonempty,
convex, closed values that is upper hemicontinuous. Then $F$ has a fixed point.
\end{theorem}

\begin{remark}
This theorem extends Kakutani to infinite dimensions, just as Schauder's theorem extends
Brouwer's.
\end{remark}

% ------------------------------------------------------------
\section{Application: Nash Equilibrium}
% ------------------------------------------------------------

\begin{definition}[Normal form game]
\index{game!normal form}
\index{Nash equilibrium}
A \emph{normal form game} with $n$ players is a triple
$\Gamma = (N, (S_i)_{i \in N}, (u_i)_{i \in N})$ where:
\begin{itemize}
  \item $N = \{1, \ldots, n\}$ is the set of players;
  \item $S_i$ is the strategy set of player $i$;
  \item $u_i\colon S_1 \times \cdots \times S_n \to \R$ is the utility function of
    player $i$.
\end{itemize}
The set of mixed strategies of player $i$ is $\Delta(S_i) = \{\sigma_i \in \R^{S_i}
: \sigma_i \geq 0, \sum_{s \in S_i} \sigma_i(s) = 1\}$.
\end{definition}

\begin{definition}[Nash equilibrium]
A mixed strategy profile $\sigma^* = (\sigma_1^*, \ldots, \sigma_n^*)$ is a
\emph{Nash equilibrium} if for every player $i$ and every mixed strategy
$\sigma_i \in \Delta(S_i)$:
\[
  u_i(\sigma^*) \geq u_i(\sigma_i, \sigma_{-i}^*)
\]
where $\sigma_{-i}^*$ denotes the strategies of the other players.
\end{definition}

\begin{definition}[Best response correspondence]
The \emph{best response correspondence} of player $i$ is
\[
  BR_i(\sigma_{-i}) = \argmax_{\sigma_i \in \Delta(S_i)} u_i(\sigma_i, \sigma_{-i}).
\]
\end{definition}

\begin{theorem}[Nash, 1950]
\label{thm:nash}
Every finite normal form game (i.e., with $S_i$ finite for all $i$) has at least one
Nash equilibrium in mixed strategies.
\end{theorem}

\begin{proof}
Let $C = \prod_{i=1}^n \Delta(S_i)$. This is a product of simplices, hence a nonempty
compact convex subset of $\R^{\sum \abs{S_i}}$.

Define the global best response correspondence:
\[
  BR(\sigma) = BR_1(\sigma_{-1}) \times \cdots \times BR_n(\sigma_{-n}).
\]
We verify Kakutani's hypotheses:

\emph{Nonemptiness and convexity}: $BR_i(\sigma_{-i})$ is nonempty (maximum of a
continuous function on a compact set) and convex (since $u_i$ is linear in $\sigma_i$:
if $\sigma_i, \sigma_i'$ both maximize $u_i(\cdot, \sigma_{-i})$, so does any convex
combination).

\emph{Upper hemicontinuity}: Berge's maximum theorem ensures that $BR$ is u.h.c., since
$u_i$ is continuous and $\Delta(S_i)$ does not depend on $\sigma_{-i}$.

By Kakutani's theorem, $BR$ has a fixed point $\sigma^*$, i.e.,
$\sigma_i^* \in BR_i(\sigma_{-i}^*)$ for all $i$. This is a Nash equilibrium.
\end{proof}

% ------------------------------------------------------------
\section{Minimax Theorems}
% ------------------------------------------------------------

\begin{theorem}[Von Neumann's minimax theorem]
\label{thm:minimax}
\index{minimax theorem}
Let $A$ be a real $m \times n$ matrix. Then:
\[
  \max_{\sigma \in \Delta_m} \min_{\tau \in \Delta_n} \sigma^T A \tau
  = \min_{\tau \in \Delta_n} \max_{\sigma \in \Delta_m} \sigma^T A \tau.
\]
\end{theorem}

\begin{proof}
The inequality $\max\min \leq \min\max$ always holds. For equality, consider the two-player
zero-sum game defined by matrix $A$. By Nash's theorem (or directly by Kakutani), this
game has an equilibrium $(\sigma^*, \tau^*)$. The value of the game is
$v = \sigma^{*T} A \tau^* = \max_\sigma \min_\tau \sigma^T A \tau =
\min_\tau \max_\sigma \sigma^T A \tau$.
\end{proof}

\begin{corollary}
Every zero-sum matrix game has a value, and both players have optimal mixed strategies.
\end{corollary}

\begin{keyformulas}[title=\textbf{Summary}]
\begin{center}
\renewcommand{\arraystretch}{1.3}
\begin{tabular}{|l|l|}
\hline
\textbf{Theorem} & \textbf{Key hypotheses} \\
\hline
Kakutani & $C$ compact convex $\subseteq \R^n$, $F(x)$ convex, u.h.c. \\
Glicksberg--Fan & $C$ compact convex in lctvs, $F(x)$ convex closed, u.h.c. \\
Nash & Finite $n$-player game \\
von Neumann & Zero-sum matrix game \\
\hline
\end{tabular}
\end{center}
\end{keyformulas}

% ------------------------------------------------------------
\section{Exercises}
% ------------------------------------------------------------

\begin{exercise}
Show that the correspondence $F\colon [0,1] \rightrightarrows [0,1]$ defined by
$F(x) = [1-x, 1]$ if $x \leq 1/2$ and $F(x) = [0, 1-x]$ if $x > 1/2$ satisfies
Kakutani's hypotheses. Find the fixed points explicitly.
\end{exercise}

\begin{exercise}
Let $C = \{x \in \R^n : \norm{x} \leq 1\}$ be the closed unit ball. Show that the
correspondence $F(x) = \{y \in C : \inner{x}{y} \geq 0\}$ satisfies Kakutani's
hypotheses and find the set of fixed points.
\end{exercise}

\begin{exercise}[Berge's maximum theorem]
\index{Berge's theorem}
Let $X, Y$ be topological spaces, $f\colon X \times Y \to \R$ continuous, and
$\Gamma\colon X \rightrightarrows Y$ a u.h.c.\ correspondence with nonempty compact
values. Show that the maximizer correspondence $M(x) = \argmax_{y \in \Gamma(x)} f(x,y)$
is u.h.c.
\end{exercise}

\begin{exercise}
Find the Nash equilibria in mixed strategies of Rock-Paper-Scissors, defined by the
payoff matrix:
\[
  A = \begin{pmatrix} 0 & -1 & 1 \\ 1 & 0 & -1 \\ -1 & 1 & 0 \end{pmatrix}.
\]
\end{exercise}

\begin{exercise}
Show that the convexity hypothesis on $F(x)$ cannot be replaced by connectedness.
\emph{Hint: consider a circle in $\R^2$.}
\end{exercise}

\begin{exercise}
Let $\Gamma$ be a two-player game where the strategy sets are compact convex subsets of
$\R^n$ and $\R^m$ respectively, and the utility functions are continuous and quasi-concave
in each player's own strategy. Prove the existence of a Nash equilibrium in pure strategies.
\end{exercise}
